\chapter{超越似然：生成对抗网络与能量模型}

\section*{概述}
\addcontentsline{toc}{section}{概述}
贯穿本书，我们已经研究了以不同数学方式表示复杂数据分布的生成模型。PPCA 与 VAE 引入了潜变量建模与近似推断；扩散模型把这一视角扩展为逐步加噪与去噪过程；基于得分的模型则进一步用学习到的得分场和随机动力学重新解释去噪。归一化流与自回归模型展示了如何借助精心设计的变换或分解，保留精确似然。

最后一章转向一个不同的问题：如果可处理的似然不再是中心对象，会发生什么？在许多重要的生成模型中，数据分布并不是由一个可以逐点计算的、显式归一化的概率密度来表示。相反，学习通过比较、几何关系或未归一化的标量景观来进行。

生成对抗网络（Generative Adversarial Networks, GAN）首先展示了这种转变。生成器学习产生难以与真实数据区分的样本，而判别器提供比较信号。Wasserstein GAN 则进一步改进了这一对抗思想：它不再使用基于分类的比较，而是采用建立在最优传输基础上的分布几何距离。

能量模型（Energy-Based Models, EBM）提供了另一条超越显式似然的路径。它不直接定义归一化密度，而是为每个配置赋予一个标量能量。低能量标记数据空间中合理的区域，高能量则抑制不合理的配置。因此，学习的任务变成塑造一个能量景观，而不是直接最大化某个可处理的似然。

本章旨在补全本书建立起来的概念版图。GAN、Wasserstein GAN 与能量模型共同表明：即使无法得到归一化概率，生成学习仍然可以进行——通过比较样本、度量几何差异，或者学习其梯度能够引导样本在数据空间中移动的标量景观。

\section*{概念图}
\addcontentsline{toc}{section}{概念图}

本章研究不要求可处理似然计算的生成模型。核心主题是：当概率是隐式的、未归一化的，或者被比较关系所取代时，学习如何继续进行。

\begin{itemize}

    \item[\(\rightarrow\)] \textbf{首先，我们把 GAN 介绍为“通过比较来学习”。}  
    生成器把潜在噪声映射为类似数据的样本，而判别器学习区分生成样本与真实数据。生成器则通过试图让这种比较失效来学习。

    \item[\(\rightarrow\)] \textbf{随后，我们从分布层面解释对抗博弈。}  
    对于固定的生成器，最优判别器估计数据分布与生成器分布之间的局部密度比。把这一判别器代回 GAN 目标函数，会得到 Jensen--Shannon 散度。

    \item[\(\rightarrow\)] \textbf{接着，我们讨论 GAN 优化及其局限。}  
    原始极小极大生成器更新在判别器非常确信时可能产生很弱的梯度。非饱和损失可以改善生成器梯度，但模式坍缩以及基于重叠的比较仍然是重要限制。

    \item[\(\rightarrow\)] \textbf{然后，我们从基于重叠的比较转向 Wasserstein 几何。}  
    Wasserstein GAN 用 $1$-Wasserstein 距离取代基于分类的差异度量，衡量把概率质量从一个分布搬运到另一个分布所需的距离。由此引出评论器（critic）、Lipschitz 约束以及具有几何意义的训练。

    \item[\(\rightarrow\)] \textbf{接下来，本章转向把能量模型理解为标量景观。}  
    能量模型为每个配置赋予标量能量 $E_\theta(x)$。低能量标记数据空间中的合理区域，高能量抑制不合理配置。模型通常定义的是未归一化密度，而不是可处理的似然。

    \item[\(\rightarrow\)] \textbf{随后，我们推导 EBM 学习中的对比结构。}  
    能量模型的最大似然会降低数据样本的能量，同时提高来自模型分布的样本能量。这暴露出核心困难：学习需要从当前模型中进行采样。

    \item[\(\rightarrow\)] \textbf{最后，我们连接能量、得分与扩散。}  
    在未归一化能量模型中，低能量对应高概率，而得分等于负能量梯度。因此，可以把扩散模型看成在多个噪声尺度上学习类似能量梯度的方向，从而提供一条稳定的多尺度路径，绕开显式似然。

\end{itemize}

\noindent 本章讨论的三条“超越似然”路线概括如下。
\begin{center}
\begin{tikzpicture}[
  >=Latex,
  every node/.style={font=\small},
  box/.style={draw,rounded corners=2pt,align=center,inner sep=6pt,
             text width=3.15cm,minimum height=1.02cm},
  flow/.style={->,thick}
]
\node[font=\normalsize\bfseries] at (6.0,2.65)
  {超越似然：概览};

\node[box] (gan) at (1.7,0.85) {
  \textbf{GAN 比较}\\[2pt]
  生成器样本\\
  $\leftrightarrow$ 判别器\\
  反馈
};
\node[box] (wgan) at (6.0,0.85) {
  \textbf{WGAN 几何}\\[2pt]
  传输距离\\
  取代\\
  分类
};
\node[box] (ebm) at (10.3,0.85) {
  \textbf{EBM 景观}\\[2pt]
  能量赋予\\
  相对\\
  合理性
};
\draw[flow] (gan) -- (wgan);
\draw[flow] (wgan) -- (ebm);

\node[box] (ratio) at (1.7,-2.15) {
  \textbf{最优}\\[-1pt]
  \textbf{判别器}\\[2pt]
  局部密度比\\
  信息
};
\node[box] (energy) at (6.0,-2.15) {
  \textbf{未归一化密度}\\[2pt]
  $p_\theta(x)\propto e^{-E_\theta(x)}$\\
  全局\\
  配分函数
};
\node[box] (score) at (10.3,-2.15) {
  \textbf{回到得分}\\[2pt]
  $\nabla_x\log p_\theta(x)$\\
  $=-\nabla_xE_\theta(x)$
};
\draw[flow] (gan) -- (ratio);
\draw[flow] (ebm) -- (energy);
\draw[flow] (energy) -- (score);
\end{tikzpicture}
\end{center}

\section{生成对抗网络：通过比较来学习}

到目前为止介绍的许多生成模型都依赖显式概率分布。无论使用精确似然还是精心构造的近似，学习通常都被表述为给观测数据赋予高概率的问题。这提供了原则清晰的基础，但也对模型如何设计与训练施加了强结构约束。

于是自然会产生一个更简单的问题：
\[
\textit{为了生成逼真的数据，我们真的必须计算概率吗？}
\]
如果目标只是生成看起来像来自数据分布的样本，那么也许通过比较进行学习就足够了。我们不再询问一个样本有多大似然，而是询问它能否与真实数据区分。这样，学习就通过匹配样本来进行，而不是通过计算密度。

生成对抗网络（GAN）~\cite{goodfellow2014generative} 正是围绕这一思想构建。它通过生成令人信服的样本这一能力来定义生成模型，而无须规定或计算概率密度。训练由比较驱动：生成样本与真实数据相互对抗，模型不断学习，使二者越来越难以区分。

\subsection{生成器作为隐式采样机制}

GAN 从一个\emph{生成器（generator）}开始，它是参数化函数
\[
x = G_\theta(z),
\qquad z \sim p(z),
\]
其中 $z$ 是从简单固定分布（如标准高斯分布或均匀分布）中抽取的潜变量。生成器 $G_\theta$ 通常由神经网络实现，把这一简单随机源映射到数据空间。

映射
\[
z \longmapsto G_\theta(z)
\]
把简单噪声转化为有结构、类似真实数据的样本。对 $z\sim p(z)$ 采样并送入 $G_\theta$，生成器会在 $x$ 上诱导出一个分布，记为 $p_\theta(x)$。关键在于，这个分布只被\emph{隐式}定义：我们可以从中采样，但通常无法计算给定点处的密度值。

这使 GAN 与前文基于似然的模型有所区别。与 VAE 不同，潜变量 $z$ 不是从数据中推断得到的，也不学习后验分布；与扩散模型不同，$z$ 并非由某个前向加噪过程产生；与归一化流或自回归模型不同，生成器不要求可逆，也无须计算 Jacobian 行列式或显式归一化 $p_\theta(x)$。因此，生成器更像一个确定性模拟器：它规定\emph{如何生成数据}，却不规定\emph{某个给定数据点有多大可能出现}。

训练完成后，从 GAN 采样非常直接。先从潜在噪声分布中抽取 $z\sim p(z)$，再计算 $x=G_\theta(z)$。生成阶段不再需要判别器；判别器只在训练时提供塑造生成器的信号。因此，GAN 能够直接采样，却不能直接计算似然。

\subsection{判别器与对抗博弈}

如果模型只会生成样本，就立即面临一个基本问题：
\[
\textit{我们如何判断这些样本是否足够好？}
\]
在基于似然的模型中，概率回答了这个问题：如果样本在模型下具有高似然，就认为它们是好的。而在 GAN 中，似然不可获得，没有一个可处理的 $p_\theta(x)$ 可供计算或最大化。

为了解决这一问题，GAN 引入第二个学习组件：\emph{判别器（discriminator）}。判别器并不测量概率密度，而是对两种样本来源进行数据驱动的比较：
\[
x\sim p_{\text{data}}(x),
\qquad
x=G_\theta(z),\quad z\sim p(z).
\]
第一种来源是真实数据样本，第二种来源是生成样本。

形式上，判别器是函数
\[
D_\phi:\mathbb{R}^d\to(0,1),
\]
由参数 $\phi$ 决定。对于输入 $x$，$D_\phi(x)$ 被解释为判别器认为 $x$ 来自真实数据分布而非生成器的置信度。因此，判别器被训练为二分类器：
\[
D_\phi(x)\approx 1
\quad\text{对真实数据，}
\qquad
D_\phi(G_\theta(z))\approx 0
\quad\text{对生成数据。}
\]

判别器本身并不是最终目的。它既不定义概率密度，也不是训练结束后我们真正想使用的对象。它的作用，是在真实样本与生成样本之间提供一个不断变化、自适应的差异度量。生成器完全通过这一反馈来学习，不断调整参数，使生成样本越来越难与真实样本区分。

这种相互作用由极小极大目标形式化：
\[
\boxed{
\min_\theta \max_\phi
\;\;
\mathbb{E}_{x \sim p_{\text{data}}}
\big[\log D_\phi(x)\big]
+
\mathbb{E}_{z \sim p(z)}
\big[\log (1 - D_\phi(G_\theta(z)))\big].
}
\]
判别器试图最大化该目标：给真实样本高分、给生成样本低分。生成器则试图最小化它，使 $D_\phi(G_\theta(z))$ 变大，从而让判别器难以把生成样本判为伪造。

\begin{tcolorbox}[
  colback=blue!4,
  colframe=blue!50!black,
  title={如何理解极小极大目标},
  boxrule=0.6pt,
  arc=2mm,
  left=2mm,
  right=2mm,
  top=1mm,
  bottom=1mm
]
表达式
\[
\min_\theta \max_\phi \; \mathcal{L}(\theta,\phi)
\]
应被理解为一个双人博弈。

在生成器参数 $\theta$ 固定时，内层最大化询问：
\[
\textit{面对当前生成器，最优判别器是什么？}
\]
这一步训练 $D_\phi$ 区分真实样本与生成样本。

随后，把判别器视为固定，外层最小化询问：
\[
\textit{生成器应如何改变，才能让真实样本与生成样本更难区分？}
\]
实践中，这两步交替执行：先改进判别器，再改进生成器。
\end{tcolorbox}

尽管这个目标看起来类似分类损失，其作用本质上是生成性的。判别器在两个分布之间定义一个学习到的比较，而生成器不断重塑自己的输出分布，使这种比较逐渐失去区分能力。

\subsection{最优判别器}

为了从分布层面理解对抗训练，考虑一个理想化情形：生成器 $G_\theta$ 固定。对 $z\sim p(z)$ 采样并令 $x=G_\theta(z)$，会诱导出隐式生成器分布 $p_\theta(x)$。在生成器固定后，判别器的任务就退化为在
\[
p_{\text{data}}(x)
\qquad\text{与}\qquad
p_\theta(x)
\]
之间进行二分类。

判别器目标可以写成
\[
\mathbb{E}_{x \sim p_{\text{data}}}\!\big[\log D(x)\big]
+
\mathbb{E}_{x \sim p_\theta}\!\big[\log(1 - D(x))\big].
\]
把期望写成积分，得到
\[
\int
p_{\text{data}}(x)\log D(x)\,dx
+
\int
p_\theta(x)\log(1-D(x))\,dx.
\]
由于目标函数对每个 $x$ 都可以分解，因此最优判别器可以逐点推导。对固定的 $x$，我们关于标量 $D(x)$ 最大化
\[
p_{\text{data}}(x)\log D(x)
+
p_\theta(x)\log(1-D(x)).
\]

对 $D(x)$ 求导，得到
\[
\frac{p_{\text{data}}(x)}{D(x)}
-
\frac{p_\theta(x)}{1-D(x)}.
\]
令导数为零，得到
\[
\frac{p_{\text{data}}(x)}{D(x)}
=
\frac{p_\theta(x)}{1-D(x)}.
\]
解出 $D(x)$，得到最优判别器
\[
\boxed{
D^*(x)
=
\frac{p_{\text{data}}(x)}
{p_{\text{data}}(x)+p_\theta(x)}.
}
\]

这一公式给出一个重要解释：即使两种密度都没有显式提供给模型，最优判别器仍然是一个局部密度比对象。图~\ref{fig:gan-optimal-discriminator-geometry} 直观展示了这一点：数据密度占优处，其输出趋近于一；生成器密度占优处，其输出趋近于零；两种密度相等处，输出恰好为二分之一。

\begin{figure}[H]
\centering
\includegraphics[width=0.98\textwidth]{figures/ch08/figure8_1_optimal_discriminator.pdf}
\caption{\textbf{最优判别器反映局部密度的相对平衡。}
对于固定生成器，$D^*(x)=p_{\mathrm{data}}(x)/(p_{\mathrm{data}}(x)+p_\theta(x))$。在数据密度占优处它趋近于一，在生成器密度占优处趋近于零，而在两种密度相等处恰为 $1/2$。}
\label{fig:gan-optimal-discriminator-geometry}
\end{figure}

\subsection{从最优判别器到 Jensen--Shannon 散度}

得到固定生成器下的最优判别器之后，我们可以进一步问：整个对抗博弈在分布层面究竟代表什么？为此，把判别器目标写成一个\emph{价值函数（value function）}，它通过判别器区分真实样本与生成样本的能力，为一个生成器--判别器组合打分。

定义
\[
V(G_\theta,D)
=
\mathbb{E}_{x \sim p_{\text{data}}}\!\big[\log D(x)\big]
+
\mathbb{E}_{x \sim p_\theta}\!\big[\log(1-D(x))\big].
\]
对于固定的 $G_\theta$，关于 $D$ 最大化 $V(G_\theta,D)$ 就得到针对该生成器的最优判别器。然后在这个最优判别器处计算价值函数，就可以识别生成器实际上在隐式最小化哪一种分布差异。

利用
\[
D^*(x)
=
\frac{p_{\text{data}}(x)}
{p_{\text{data}}(x)+p_\theta(x)},
\]
经过代数化简，得到
\[
V(G_\theta,D^*)
=
-\log 4
+
2\,\mathrm{JS}\!\left(p_{\text{data}}\,\|\,p_\theta\right),
\]
其中 $\mathrm{JS}(\cdot\|\cdot)$ 表示 Jensen--Shannon 散度。

常数项 $-\log 4$ 与生成器无关，因此不影响优化。所以，在判别器达到最优的理想化情形下，GAN 训练等价于最小化数据分布与生成器分布之间的 Jensen--Shannon 散度。

上式引入了概率分布之间一种新的散度：\emph{Jensen--Shannon 散度}。要理解它，可以回忆 VAE 章节在证据下界中出现过的 Kullback--Leibler（KL）散度。KL 散度按如下方式比较两个分布：
\[
\mathrm{KL}(p\,\|\,q)
=
\int p(x)\log\frac{p(x)}{q(x)}\,dx.
\]
它衡量用 $q$ 近似 $p$ 所付出的代价。然而，KL 散度是不对称的：
\[
\mathrm{KL}(p\,\|\,q)
\neq
\mathrm{KL}(q\,\|\,p),
\]
并且对支撑集不匹配十分敏感。如果在某一区域 $p(x)>0$ 而 $q(x)=0$，那么 $\mathrm{KL}(p\,\|\,q)$ 会变为无穷大。

Jensen--Shannon 散度通过让两个分布都与一个共享的中间分布进行比较来修改这一方式。给定两个分布 $p$ 与 $q$，定义它们的混合分布
\[
m(x)
=
\frac{1}{2}\big(p(x)+q(x)\big).
\]
可以把这个混合分布理解为：先以相同概率在 $p$ 与 $q$ 之间选择一个分布，再从被选中的分布中采样。

Jensen--Shannon 散度定义为
\[
\boxed{
\mathrm{JS}(p\,\|\,q)
=
\frac{1}{2}\mathrm{KL}(p\,\|\,m)
+
\frac{1}{2}\mathrm{KL}(q\,\|\,m).
}
\]
因此，它不要求一个分布直接解释另一个分布，而是询问两个分布分别与它们共同的混合分布相距多远。这使其具有对称性：
\[
\mathrm{JS}(p\,\|\,q)
=
\mathrm{JS}(q\,\|\,p),
\]
并且有界：
\[
0
\leq
\mathrm{JS}(p\,\|\,q)
\leq
\log 2.
\]

这正是把最优判别器代回原始 GAN 价值函数后出现的散度。因此，在判别器最优的理想化条件下，生成器并不是在最大化似然；它是在调整自己的隐式分布 $p_\theta$，以减小
\[
\mathrm{JS}\!\left(p_{\text{data}}\,\|\,p_\theta\right).
\]

这揭示了对抗训练在分布层面的含义。GAN 并不构造显式密度，而是执行\emph{通过比较进行的分布匹配}。判别器估计真实分布与生成分布有多容易区分，生成器则通过降低这种可区分性来学习。

全局最优点满足
\[
p_\theta(x)=p_{\text{data}}(x).
\]
此时，处处都有
\[
D^*(x)=\frac{1}{2},
\]
因为真实样本与生成样本已经无法区分。

\section{几何优化与 Wasserstein GAN}

前面的分析表明，在最优判别器下，原始 GAN 目标对应于最小化数据分布与生成器隐式分布之间的 Jensen--Shannon 散度。这给出了清晰的理论解释，同时也暴露出一个局限：当 $p_{\text{data}}$ 与 $p_\theta$ 很少甚至完全不重叠时，JS 散度可能进入饱和状态。此时判别器几乎可以完美地区分真实样本与生成样本，却几乎不能告诉生成器应当如何向数据分布移动。这促成了一个关键转变：从基于分布重叠的比较，转向几何距离。Wasserstein GAN~\cite{arjovsky2017wasserstein} 不再询问两个分布“是否可区分”，而是询问：为了把一个分布变成另一个分布，需要把概率质量移动多远。

\subsection{生成器优化与非饱和损失}

在转向 Wasserstein GAN 之前，我们先讨论原始 GAN 框架中的一种常用实践改进：改变生成器损失，以改善训练时的梯度信号。
原始 GAN 目标写作极小极大问题
\[
\min_\theta \max_\phi
\;
\mathbb{E}_{x \sim p_{\text{data}}}\big[\log D_\phi(x)\big]
+
\mathbb{E}_{z \sim p(z)}\big[\log (1 - D_\phi(G_\theta(z)))\big].
\]
对判别器而言，其解释很直接：真实样本得分越高、生成样本得分越低，奖励越大。但对生成器而言，情况更微妙。

若固定判别器，并关于生成器参数 $\theta$ 最小化同一个价值函数，唯一依赖 $\theta$ 的项是
\[
\mathbb{E}_{z \sim p(z)}
\big[\log (1 - D_\phi(G_\theta(z)))\big].
\]
这一目标要求生成器降低判别器认为生成样本是伪造数据的置信度。

困难常出现在训练早期。此时判别器通常可以非常自信地识别生成样本：
\[
D_\phi(G_\theta(z))\approx 0.
\]
于是
\[
1-D_\phi(G_\theta(z))\approx 1,
\qquad
\log(1-D_\phi(G_\theta(z)))\approx 0.
\]
生成器目标在这一区域因此接近平坦，即使生成样本仍然很差，反向传给生成器的梯度也可能非常小。

因此，GAN 实践中常采用另一个生成器目标，称为\emph{非饱和损失（non-saturating loss）}：
\[
\min_\theta
\;
\mathbb{E}_{z \sim p(z)}
\big[
-\log D_\phi(G_\theta(z))
\big].
\]
这个损失并不是对前一个表达式做代数化简得到的，而是把生成器的目标从降低判别器的“伪造”置信度，改为提高其“真实”置信度。

两种目标不同，但期望达到的平衡点相同：生成样本最终应与真实数据无法区分。它们真正的差异在于梯度行为。当
\[
D_\phi(G_\theta(z))\approx 0
\]
时，非饱和项
\[
-\log D_\phi(G_\theta(z))
\]
很大，因此生成器能获得更强的学习信号。这就是实践中通常更偏好非饱和损失的原因，尽管原始极小极大目标与 Jensen--Shannon 散度之间的理论联系更加直接。

这种优化行为有助于理解 GAN 的优势与弱点。一方面，判别器会惩罚让生成样本显得不真实的特征，因此对抗训练能够鼓励清晰、视觉上更可信的输出，而不是模糊的平均结果。另一方面，同一机制也可能导致\emph{模式坍缩（mode collapse）}。由于生成器的目标是欺骗判别器，它可能只把数据分布中的一部分模式建模得很好。只要这一子集足以迷惑判别器，其他模式就可能被忽略。

基于似然的模型如果不能覆盖观测到的模式，会直接受到惩罚，因为似然要求模型给数据分配概率质量。相比之下，GAN 优化的是可区分性，而不是显式的概率覆盖。因此，非饱和损失虽然改善了生成器梯度，却没有改变对抗比较的底层几何。当两个分布几乎不重叠时，基于分类的差异度量可能只告诉我们它们“可区分”，却没有告诉我们一个分布应该朝哪里移动才能接近另一个分布。这就需要一种更具几何意义的分布差异度量。

\subsection{从重叠关系到 Wasserstein 几何}

为了发展这一思想，可以从几何角度理解概率分布。不要把分布只看成静态密度，而把其概率质量想象成可以在空间中搬运的物质。如果把一个分布变成另一个分布只需要做很小的质量位移，那么两个分布应被认为很接近；如果必须把大量质量搬运很远，它们就应被认为相距很远。

Wasserstein 距离把这种直觉形式化。在 Wasserstein GAN 中，相关的是 $1$-Wasserstein 距离，也称 Earth Mover's distance（推土机距离）。它衡量把概率质量从一个分布搬运到另一个分布的最小代价，其中代价与搬运距离成正比。

一个关键理论结果——Kantorovich--Rubinstein 对偶性——给出了这种距离的实用表达。对于 $\mathbb{R}^d$ 上两个概率分布 $p$ 与 $q$，$1$-Wasserstein 距离可写为
\[
\boxed{
W(p,q)
=
\sup_{\|f\|_{L} \le 1}
\left(
\mathbb{E}_{x \sim p}[f(x)]
-
\mathbb{E}_{x \sim q}[f(x)]
\right).
}
\]
这里，$f$ 是实值评分函数，而约束
\[
\|f\|_{L}\leq 1
\]
表示 $f$ 必须是 $1$-Lipschitz 的：其函数值相对于输入空间中的距离，变化速度不能超过单位速率。

表达式
\[
\mathbb{E}_{x \sim p}[f(x)]
-
\mathbb{E}_{x \sim q}[f(x)]
\]
衡量该评分函数在平均意义下把两个分布区分得有多强。上确界意味着：在满足 Lipschitz 约束的所有评分函数中，选择区分能力最强的那个。

这一约束至关重要。没有它，评分函数可以给一个分布中的样本任意大的值、给另一个分布中的样本任意小的值，从而使比较失去意义。Lipschitz 约束迫使评分差异反映数据空间的几何结构：只有当对应点彼此相距较远时，才允许出现很大的评分差异。

需要注意，这里的“score”不同于基于得分的扩散模型中的得分。在后者中，得分指 $\nabla_x\log p(x)$，即对数密度的梯度；而在 WGAN 中，评论器输出 $f_\phi(x)$ 是用于比较的实值标量，而不是概率。对评论器施加的是 Lipschitz 行为约束，而非概率归一化约束。

\begin{tcolorbox}[
  colback=blue!4,
  colframe=blue!50!black,
  title={Wasserstein 距离的一个计算示例},
  boxrule=0.6pt,
  arc=2mm,
  left=2mm,
  right=2mm,
  top=1mm,
  bottom=1mm
]
考虑一维例子：$p$ 把全部概率质量放在 $x=0$，而 $q$ 把全部概率质量放在 $x=10$。从 $p$ 采样总会得到 $0$，从 $q$ 采样总会得到 $10$，于是
\[
\mathbb{E}_{x\sim p}[f(x)] = f(0),
\qquad
\mathbb{E}_{x\sim q}[f(x)] = f(10),
\qquad
W(p,q)
=
\sup_{\|f\|_L\leq 1}
\big(f(0)-f(10)\big).
\]

约束 $\|f\|_L\leq 1$ 意味着
\[
|f(x_1)-f(x_2)|\leq |x_1-x_2|
\quad
\text{对任意 }x_1,x_2.
\]
因此，每个合法的 $f$ 都满足
\[
f(0)-f(10)\leq |0-10|=10.
\]
函数 $f(x)=-x$ 达到这一上界，因为
\[
f(0)-f(10)=0-(-10)=10.
\]
所以 $W(p,q)=10$，与搬运直觉完全一致：全部概率质量都必须移动距离 $10$。
\end{tcolorbox}

这个例子说明了为什么在分布完全不重叠时，Wasserstein 距离仍然具有信息量。基于重叠的散度能够检测支撑集彼此分离，却未必能反映两个分布相距一个单位还是十个单位。相比之下，Wasserstein 距离会直接度量空间分离程度。图~\ref{fig:gan-js-wasserstein-geometry} 对一族平移高斯分布展示了相同差异：随着重叠消失，JS 散度很快接近饱和，而 Wasserstein 距离仍持续追踪空间上的分离程度。

\begin{figure}[H]
\centering
\begin{minipage}[t]{0.48\textwidth}
  \vspace{0pt}\centering
  \includegraphics[width=\linewidth]{figures/ch08/panels/figure8_2_a.pdf}
\end{minipage}\hfill
\begin{minipage}[t]{0.48\textwidth}
  \vspace{0pt}\centering
  \includegraphics[width=\linewidth]{figures/ch08/panels/figure8_2_b.pdf}
\end{minipage}

\vspace{0.8em}

\begin{minipage}[t]{0.48\textwidth}
  \vspace{0pt}\centering
  \includegraphics[width=\linewidth]{figures/ch08/panels/figure8_2_c.pdf}
\end{minipage}\hfill
\begin{minipage}[t]{0.48\textwidth}
  \vspace{0pt}
  \caption{\textbf{JS 散度会饱和，而 Wasserstein 距离仍保留几何信息。}
  各子图比较方差相同、均值逐渐分离的高斯分布。随着重叠消失，JS 散度逼近上界 $\log 2$，对进一步分离逐渐不敏感。对于这些平移分布，$1$-Wasserstein 距离为 $W_1=\Delta$，因此仍持续度量概率质量需要搬运多远。}
  \label{fig:gan-js-wasserstein-geometry}
\end{minipage}
\end{figure}


\subsection{Wasserstein GAN 目标}

Wasserstein 距离的对偶表征直接导出一个可实践的对抗目标。如果生成器诱导出分布 $p_\theta$，那么数据分布与生成器分布之间的 Wasserstein 距离可以写成
\[
W(p_{\text{data}},p_\theta)
=
\sup_{\|f\|_{L}\leq 1}
\left(
\mathbb{E}_{x\sim p_{\text{data}}}[f(x)]
-
\mathbb{E}_{x\sim p_\theta}[f(x)]
\right).
\]
由于 $p_\theta$ 中的样本通过
\[
x=G_\theta(z),
\qquad z\sim p(z)
\]
生成，因此可以写成
\[
\mathbb{E}_{x\sim p_\theta}[f(x)]
=
\mathbb{E}_{z\sim p(z)}[f(G_\theta(z))].
\]

Wasserstein GAN 用神经网络 $f_\phi$ 表示评分函数 $f$，并优化
\[
\boxed{
\min_\theta \max_{\phi \in \mathcal{L}}
\;
\mathbb{E}_{x \sim p_{\text{data}}}[f_\phi(x)]
-
\mathbb{E}_{z \sim p(z)}[f_\phi(G_\theta(z))],
}
\]
其中 $\mathcal{L}$ 表示所有使 $f_\phi$ 满足所需 Lipschitz 约束的参数集合。

关于 $\phi$ 的最大化试图找到一个评分函数：在尊重数据空间几何结构的前提下，尽可能强地区分真实样本与生成样本。随后，关于 $\theta$ 的最小化更新生成器，使这种分离越来越小。

到这里，术语也随之变化。在原始 GAN 中，负责评价样本的网络称为\emph{判别器}，其输出
\[
D_\phi(x)\in(0,1)
\]
被解释为 $x$ 是真实样本的概率。而在 Wasserstein GAN 中，对应的网络称为\emph{评论器（critic）}。它输出实值分数
\[
f_\phi(x)\in\mathbb{R},
\]
而不是概率。这里没有 sigmoid 输出，不作二分类解释，也不要求分数归一化；真正有意义的是评论器分数之间的差异。

从这个角度看，评论器是 Wasserstein 距离中最优评分函数的神经网络近似。它的目的不是作出一个硬性的真假判断，而是赋予样本分数，使其平均差异反映两个分布之间的几何分离程度。在分数不能跨数据空间任意快速变化的约束下，真实样本的平均得分应高于生成样本。由于 Wasserstein 距离会随着生成器分布的移动而连续变化，即使 $p_\theta$ 与 $p_{\text{data}}$ 没有重叠，生成器仍能获得有用梯度；因此，优化由几何邻近性而非仅由重叠程度来引导。

与原始 GAN 一样，评论器只用于训练。生成器训练完成后，采样仍然只是抽取 $z\sim p(z)$ 并返回 $G_\theta(z)$。

\subsection{施加 Lipschitz 约束}

剩下的实际问题是：训练期间如何对评论器施加 Lipschitz 约束？回忆这一约束的作用：让评论器分数以受控速率变化，使分数差异反映几何距离，而不是任意缩放。普通神经网络并不会自动满足这一性质。

最初的 WGAN 通过\emph{权重裁剪（weight clipping）}施加该约束，把评论器的每个参数限制在固定区间内。这种方法简单，却是一种非常粗糙的 Lipschitz 常数控制方式，可能削弱评论器容量，并导致优化不稳定。

更有效的替代方法是 WGAN-GP~\cite{gulrajani2017improved} 引入的\emph{梯度惩罚（gradient penalty）}。它不直接限制参数，而是限制评论器的输入--输出行为。其核心思想是控制评论器分数对输入微小变化的敏感程度。

在评论器训练目标中加入如下惩罚项：
\[
\lambda
\,
\mathbb{E}_{\hat{x}}
\left[
\big(\|\nabla_{\hat{x}} f_\phi(\hat{x})\|_2-1\big)^2
\right],
\]
其中 $\hat{x}$ 从真实样本与生成样本之间的直线插值点上采样。

梯度
\[
\nabla_{\hat{x}} f_\phi(\hat{x})
\]
衡量评论器分数对输入无穷小变化的敏感程度。因此，其欧氏范数就是评论器在 $\hat{x}$ 处的局部变化率，也就是局部斜率。惩罚 $\|\nabla_{\hat{x}} f_\phi(\hat{x})\|_2$ 偏离一的程度，会鼓励评论器在最相关的比较区域中平滑、受控地改变分数。

插值点 $\hat{x}$ 位于真实样本与生成样本之间，而这里恰好是评论器必须学习如何比较两个分布的区域。沿这些路径施加约束，既能稳定训练，又能保留 Wasserstein 距离的几何解释。系数 $\lambda$ 控制惩罚强度：较大的值会更强地施加约束，较小的值则允许更大灵活性。

通过这种方式施加 Lipschitz 约束，带梯度惩罚的 WGAN 在保留 Wasserstein 距离最优传输含义的同时，获得了更可靠的训练动力学。更广泛地说，它们说明生成学习不必依赖显式似然或归一化密度，也可以由一个学习得到的实值函数来引导；该函数的几何性质会把生成样本推向数据分布。

这自然形成通往能量模型的桥梁。与 WGAN 类似，能量模型也利用实值函数比较、排序或引导样本。但它不通过对抗博弈来表述这一思想，而是通过未归一化的能量景观定义分布。下面转向这一视角。

\section{能量模型：基础与学习}

前面几节讨论了通过比较进行学习的生成模型。在 GAN 中，判别器提供一个学习得到的测试，用来区分真实样本与生成样本；在 Wasserstein GAN 中，这种比较被几何评论器所取代，其函数值反映把概率质量搬运多远才能完成分布变换。两种方法都绕开显式似然计算，转而依赖样本之间的相对评价进行学习。

能量模型~\cite{lecun2006tutorial} 则把这一思想进一步推进。它不从生成器、判别器或显式似然出发，而是从数据空间上的单个标量函数开始：\emph{能量函数（energy function）}。该函数为每个配置 $x$ 赋予一个值
\[
E_\theta(x)\in\mathbb{R}.
\]
较低能量表示该配置与数据结构更相容，较高能量表示该配置不太合理。

能量函数有意不施加太多限制：它无须为正、无须有界，也无须归一化，而且它本身并不定义概率分布。因此，在这一阶段不能把能量直接解释为概率。更合适的理解是：它是一种偏好、代价或相容性度量。类似真实数据的配置应位于低能量区域，而不真实的配置则应位于高能量区域。

这给出了一种几何化的建模视角。能量模型在数据空间上定义一个景观：
\[
\text{低能量}
\quad\Longleftrightarrow\quad
\text{合理配置},
\qquad
\text{高能量}
\quad\Longleftrightarrow\quad
\text{不合理配置}.
\]
学习意味着塑造这一景观，使观测数据处于能量谷底，而远离数据的区域被赋予较高能量。

\subsection{作为建模基本对象的能量}

能量模型中的核心建模对象是标量能量函数
\[
E_\theta:\mathbb{R}^d\to\mathbb{R}.
\]
与归一化流或自回归模型不同，它不需要给出精确密度；与 VAE 不同，它不要求显式潜变量后验；与 GAN 不同，它也不需要在对抗博弈中训练一个独立判别器。模型直接由一个标量景观规定。

只有相对能量差异真正重要。如果给能量函数加上常数 $c$，
\[
\tilde{E}_\theta(x)=E_\theta(x)+c,
\]
所有配置的能量都移动同样的幅度，从低到高的相对排序保持不变。由于模型行为依赖“哪些配置相对于其他配置更受偏好”，这种整体常数平移没有实质影响。

这种相对性十分重要。能量函数并不是试图为每个点赋予绝对概率值，而是在编码偏好：哪些配置应当被鼓励，哪些应当被抑制。从这个意义上说，能量建模提出的是一个不同于显式密度建模的问题：
\[
\textit{模型应当把低能量区域放在数据空间的什么位置？}
\]
答案由数据学习得到。

\subsection{从能量到未归一化分布}

尽管能量函数本身不是概率分布，但可以通过 Boltzmann 形式把它与概率联系起来：
\[
p_\theta(x)
=
\frac{1}{Z_\theta}
\exp\big(-E_\theta(x)\big),
\]
其中
\[
Z_\theta
=
\int
\exp\big(-E_\theta(x)\big)\,dx
\]
称为归一化常数或\emph{配分函数（partition function）}。

负指数把低能量转化为较大的正权重，而配分函数执行一次全局归一化。图~\ref{fig:ebm-energy-to-density} 展示了从标量能量到未归一化权重，再到归一化密度的映射；它也清楚说明了为什么给所有能量加同一个常数不会改变最终密度。

\begin{figure}[H]
\centering
\begin{minipage}[t]{0.48\textwidth}
  \vspace{0pt}\centering
  \includegraphics[width=\linewidth]{figures/ch08/panels/figure8_3_a.pdf}
\end{minipage}\hfill
\begin{minipage}[t]{0.48\textwidth}
  \vspace{0pt}\centering
  \includegraphics[width=\linewidth]{figures/ch08/panels/figure8_3_b.pdf}
\end{minipage}

\vspace{0.8em}

\begin{minipage}[t]{0.48\textwidth}
  \vspace{0pt}\centering
  \includegraphics[width=\linewidth]{figures/ch08/panels/figure8_3_c.pdf}
\end{minipage}\hfill
\begin{minipage}[t]{0.48\textwidth}
  \vspace{0pt}
  \caption{\textbf{从能量到归一化密度。}
  低能量区域获得较大的未归一化权重 $\exp(-E_\theta(x))$；除以 $Z_\theta$ 后即可变成概率密度而不改变其形状。给所有能量加同一个常数，会按相同比例缩放权重与 $Z_\theta$，因此归一化后的密度保持不变。}
  \label{fig:ebm-energy-to-density}
\end{minipage}
\end{figure}


困难在于，高维空间中的 $Z_\theta$ 通常不可处理。计算它要求对所有可能配置进行积分：
\[
Z_\theta
=
\int
\exp\big(-E_\theta(x)\big)\,dx.
\]
对于图像、音频、文本表示或其他高维数据，这个积分一般不可能精确求得。

这正是能量建模的核心权衡。由于避免了精确密度模型所需的结构限制，能量函数可以非常灵活；但这种灵活性以归一化困难为代价。模型通常只定义\emph{未归一化}密度
\[
p_\theta(x)
\propto
\exp\big(-E_\theta(x)\big),
\]
而不是可处理的归一化似然。

这使能量模型不同于上一章讨论的精确密度模型。归一化流与自回归模型经过精心设计，使 $p_\theta(x)$ 可以被精确计算；能量模型则把重点放在能量景观本身。即使归一化常数未知，相对能量差异和能量梯度仍然具有实际用途。

\subsection{学习能量景观}

能量模型中的学习，就是塑造能量函数，使其反映数据结构。假设我们观察到样本
\[
x\sim p_{\text{data}}(x).
\]
理想的能量函数应给这些观测样本较低能量：
\[
E_\theta(x)
\quad\text{较小，当}\quad
x\sim p_{\text{data}}.
\]
但仅满足这一条件还不够。如果我们只是把所有位置的能量都降低，模型就无法区分类似数据的配置与不合理配置。因此，学习必须是对比性的：既要降低数据样本的能量，也要提高其他位置的能量。

非正式地说，理想行为是
\[
E_\theta(x)
\quad\text{在数据上较低},
\qquad
E_\theta(x)
\quad\text{远离数据时较高}.
\]
正是这种对比结构赋予能量学习以意义。模型不仅要奖励观测数据，还必须学习哪些备选配置应该被抑制。

理解这一点的一个有用工具是能量梯度
\[
\nabla_x E_\theta(x).
\]
它描述能量在数据空间局部如何变化。沿
\[
-\nabla_x E_\theta(x)
\]
方向移动会降低能量，从而把一个配置引向模型认为更合理的区域。因此，能量景观不仅定义标量偏好值，还定义局部的改进方向。

这把能量模型与前文 GAN 和基于得分模型中的几何视角联系起来。在 Wasserstein GAN 中，评论器赋予实值分数，分数差异用于引导生成器；在能量模型中，能量函数赋予标量值，其梯度定义配置应当如何在数据空间中移动。两者相关但并不相同：WGAN 评论器为了区分两个分布而训练，而能量函数的目标是定义一个合理性的景观。

\subsection{最大似然与对比梯度}

尽管能量模型通常作为未归一化模型来讨论，Boltzmann 形式仍然给出了一个原则明确的学习目标。如果
\[
p_\theta(x)
=
\frac{1}{Z_\theta}\exp(-E_\theta(x)),
\]
那么最大似然试图最大化
\[
\mathbb{E}_{x\sim p_{\text{data}}}
\big[\log p_\theta(x)\big].
\]
代入能量形式的密度，得到
\[
\log p_\theta(x)
=
-E_\theta(x)-\log Z_\theta,
\]
因此
\[
\mathbb{E}_{x\sim p_{\text{data}}}
\big[\log p_\theta(x)\big]
=
-\mathbb{E}_{x\sim p_{\text{data}}}
\big[E_\theta(x)\big]
-
\log Z_\theta.
\]

第一项的作用很直观：最大化它会降低数据样本的能量。第二项则阻止“到处降低能量”这一平凡解，因为给过大的空间区域分配低能量会使配分函数增大。因此，最大似然平衡两股力量：
\[
\text{降低数据处的能量}
\qquad
\text{与}
\qquad
\text{控制其他区域的能量}.
\]

关于 $\theta$ 对目标函数求导，这一点会更加清楚。数据项的梯度为
\[
-\nabla_\theta
\mathbb{E}_{x\sim p_{\text{data}}}
\big[E_\theta(x)\big]
=
-
\mathbb{E}_{x\sim p_{\text{data}}}
\big[\nabla_\theta E_\theta(x)\big].
\]
对配分函数项，使用
\[
Z_\theta
=
\int
\exp(-E_\theta(x))\,dx.
\]
求导得到
\[
\begin{aligned}
\nabla_\theta \log Z_\theta
&=
\frac{1}{Z_\theta}
\nabla_\theta
\int
\exp(-E_\theta(x))\,dx \\[4pt]
&=
\frac{1}{Z_\theta}
\int
\exp(-E_\theta(x))
\big(-\nabla_\theta E_\theta(x)\big)\,dx \\[4pt]
&=
-
\mathbb{E}_{x\sim p_\theta}
\big[\nabla_\theta E_\theta(x)\big].
\end{aligned}
\]
由于对数似然中包含 $-\log Z_\theta$，其梯度贡献为
\[
-\nabla_\theta \log Z_\theta
=
\mathbb{E}_{x\sim p_\theta}
\big[\nabla_\theta E_\theta(x)\big].
\]

合并两项，最大似然梯度为
\[
\boxed{
\nabla_\theta
\mathbb{E}_{x\sim p_{\text{data}}}
\big[\log p_\theta(x)\big]
=
-
\mathbb{E}_{x\sim p_{\text{data}}}
\big[\nabla_\theta E_\theta(x)\big]
+
\mathbb{E}_{x\sim p_\theta}
\big[\nabla_\theta E_\theta(x)\big].
}
\]

这一表达式是能量学习的核心。它包含两个部分：
\[
\text{数据项}
\quad
-
\mathbb{E}_{p_{\text{data}}}
\big[\nabla_\theta E_\theta(x)\big],
\qquad
\text{模型项}
\quad
+
\mathbb{E}_{p_\theta}
\big[\nabla_\theta E_\theta(x)\big].
\]

数据项降低观测样本的能量，模型项提高从当前模型分布抽取的样本能量。这通常被描述为\emph{正相位（positive phase）}与\emph{负相位（negative phase）}之间的对比：正相位使数据更受偏好，而负相位则对模型当前分配了过多概率质量的配置施加反向压力。

这一对比梯度暴露出能量模型最核心的实际困难。对数据分布的期望可以直接用训练集近似；而对模型分布的期望
\[
\mathbb{E}_{x\sim p_\theta}
\big[\nabla_\theta E_\theta(x)\big]
\]
要困难得多，因为从 $p_\theta(x)\propto \exp(-E_\theta(x))$ 本身采样就是一个非平凡问题。

这使学习与采样耦合在一起。为了改进能量函数，我们需要当前模型的样本；但要生成好的模型样本，又需要一个已经被良好塑造的能量景观。这种循环依赖，是能量模型虽然强大却难以训练的主要原因之一。

下一节将更明确地展开这一联系。由于能量模型本身不提供直接采样器，生成通常被表述为在能量景观上的移动。这自然导向能量梯度动力学、Langevin 采样，以及能量建模与前文基于得分扩散方法之间的联系。

\section{能量梯度、得分场与扩散}

上一节表明，能量模型通过未归一化的能量景观间接定义概率分布：
\[
p_\theta(x)
\propto
\exp(-E_\theta(x)).
\]
这带来很强的建模灵活性，却也产生一个实际困难：模型本身不提供直接采样机制。为了生成样本，我们必须在能量景观中移动。

这正是能量梯度成为核心对象的地方。标量能量 $E_\theta(x)$ 告诉我们一个配置有多合理，而其梯度
\[
\nabla_x E_\theta(x)
\]
则告诉我们能量在数据空间局部如何变化。沿
\[
-\nabla_x E_\theta(x)
\]
方向移动会降低能量，因此把配置推向模型认为更加合理的区域。从这个意义上说，负能量梯度就是一种局部驱动力。

本节把这一思想与本书前面建立的基于得分方法连接起来。核心结论简单但十分重要：对于形如 $p(x)\propto \exp(-E(x))$ 的分布，其得分恰好就是负能量梯度：
\[
\nabla_x\log p(x)
=
-\nabla_x E(x).
\]
因此，能量模型与基于得分模型关系非常紧密。二者真正的差异，与其说是所使用的局部方向不同，不如说在于该方向如何参数化、如何学习以及如何用于生成。

\subsection{把采样看作能量引导的动力学}

能量模型给合理配置较低能量，但并不会直接告诉我们如何从对应分布中抽取相互独立的样本。与归一化流不同，它不存在从基础分布到数据空间的可逆映射；与自回归模型不同，它也没有可以从左到右顺序采样的条件分布乘积。因此，采样通常被表述为一个迭代动力学过程。

基本思想是先初始化一个配置 $x_0$，然后反复更新，使其向低能量区域移动。纯确定性的更新可以写为
\[
x_{t+1}
=
x_t
-
\eta\nabla_x E_\theta(x_t),
\]
其中 $\eta>0$ 是步长。这很像梯度下降，但此时“下降”发生在数据空间中，而不是参数空间中；被移动的是配置 $x_t$ 本身，使其走向更低能量。

然而，仅有确定性下降还不足以构成采样。它可能陷入能量景观的某个局部极小点，而无法探索多个合理区域。为了得到采样过程，需要加入随机性：
\[
\boxed{
x_{t+1}
=
x_t
-
\eta \nabla_x E_\theta(x_t)
+
\sqrt{2\eta}\,\xi_t,
\qquad
\xi_t\sim\mathcal{N}(0,I).
}
\]
这就是与能量函数相对应的离散 Langevin 型更新。

确定性项
\[
-\eta\nabla_x E_\theta(x_t)
\]
推动当前配置走向低能量，而噪声项
\[
\sqrt{2\eta}\,\xi_t
\]
注入随机性，防止动力学简单坍缩到单一极小点，并使其能够探索更广泛的能量景观。在适当条件与合适的连续时间解释下，这种“梯度加噪声”的动力学可以以能量模型分布为平稳分布。

这正是前面基于得分章节研究过的 Langevin 思想。改变的只是漂移场的写法：在那里，漂移使用得分 $\nabla_x\log p(x)$ 表示；这里则使用负能量梯度 $-\nabla_xE_\theta(x)$ 表示。下一小节会明确给出二者的联系。

\subsection{能量梯度与得分}

回忆密度 $p(x)$ 的得分定义为
\[
s(x)
=
\nabla_x\log p(x).
\]
对于能量模型，
\[
p_\theta(x)
=
\frac{1}{Z_\theta}\exp(-E_\theta(x)).
\]
取对数得到
\[
\log p_\theta(x)
=
-E_\theta(x)-\log Z_\theta.
\]
现在对 $x$ 求导。由于配分函数 $Z_\theta$ 依赖 $\theta$，但不依赖 $x$，所以
\[
\nabla_x\log Z_\theta=0.
\]
因此
\[
\boxed{
\nabla_x\log p_\theta(x)
=
-\nabla_xE_\theta(x).
}
\]

这一恒等式是能量建模与基于得分建模之间的桥梁。能量分布的得分场恰好等于能量函数的负梯度。因此，朝高概率区域移动的局部方向，与朝低能量区域移动的方向完全相同。

区别在于对象如何表示。在经典能量模型中，我们参数化一个标量函数 $E_\theta(x)$，再通过求导得到驱动向量场：
\[
x
\quad\longmapsto\quad
-\nabla_xE_\theta(x).
\]
而在基于得分的模型中，通常直接参数化向量场：
\[
s_\theta(x)
\approx
\nabla_x\log p(x).
\]
因此，能量模型从标量景观出发并导出向量场，而得分模型往往直接学习向量场本身。

这一差别很有意义。一个标量能量天然会产生保守梯度场：向量场来源于同一个底层标量函数的梯度。直接学习的得分网络作为向量场更加灵活，但除非额外施加结构，它未必被显式表示为某个标量能量的梯度。不过，就生成建模而言，二者在概念上的联系依然很强：它们都在学习把样本引向更高数据密度区域的局部方向。

\subsection{把得分匹配理解为能量梯度学习}

上一节表明，能量模型的最大似然会产生一个包含模型分布样本的对比梯度。这造成实际困难，因为从模型本身采样就很难。得分匹配（score matching）~\cite{hyvarinen2005estimation} 提供了另一条路径：不去匹配归一化概率，而是匹配局部得分场。

在基于得分的章节中，我们已经利用 Fisher 散度和分部积分详细推导了得分匹配。这里仅回顾理解能量视角所需的结果。目标是学习模型得分
\[
s_\theta(x)
\approx
\nabla_x\log p_{\text{data}}(x).
\]
对于能量模型，
\[
s_\theta(x)
=
\nabla_x\log p_\theta(x)
=
-\nabla_xE_\theta(x).
\]
因此，可以把得分匹配解释为学习一个能量景观，使其负梯度与数据得分一致。

忽略与 $\theta$ 无关的常数项，经典得分匹配目标可以写为
\[
J_{\mathrm{SM}}(\theta)
=
\mathbb{E}_{x\sim p_{\text{data}}}
\left[
\frac{1}{2}\|s_\theta(x)\|_2^2
+
\nabla_x\cdot s_\theta(x)
\right].
\]
代入
\[
s_\theta(x)=-\nabla_xE_\theta(x)
\]
得到能量形式的目标：
\[
\boxed{
J_{\mathrm{SM}}(\theta)
=
\mathbb{E}_{x\sim p_{\text{data}}}
\left[
\frac{1}{2}
\|\nabla_xE_\theta(x)\|_2^2
-
\nabla_x^2E_\theta(x)
\right],
}
\]
其中
\[
\nabla_x^2E_\theta(x)
=
\sum_{j=1}^d
\frac{\partial^2E_\theta(x)}{\partial x_j^2}
\]
是能量函数的 Laplacian。

这里最重要的不是具体推导——前文已经完成了这些细节——而是“被匹配的对象”发生了变化。最大似然比较归一化概率分布，并会产生一个模型期望项；得分匹配比较对数密度的局部梯度，并绕开配分函数，因为
\[
\nabla_x\log Z_\theta=0.
\]
用能量的语言说，它通过导数学习数据样本附近能量景观的形状，而无须精确计算归一化密度。

这就是为什么从能量视角看，得分匹配如此自然。它不询问某个数据点的概率值，而是询问：概率应当沿哪个局部方向增大？等价地，能量应当沿哪个方向减小？

\subsection{把扩散理解为多尺度能量梯度学习}

现在可以清楚表述能量模型与基于得分扩散之间的关系。能量模型定义标量景观 $E_\theta(x)$。如果该景观对应未归一化密度
\[
p_\theta(x)\propto \exp(-E_\theta(x)),
\]
那么其得分为
\[
\nabla_x\log p_\theta(x)
=
-\nabla_xE_\theta(x).
\]
因此，Langevin 采样中使用的局部方向，与基于得分模型试图直接学习的对象完全相同：指向更高概率、也就是更低能量的方向。

从这一视角出发，可以把基于得分的扩散模型~\cite{song2019generative,song2021scorebased} 看作在一系列含噪分布上学习随时间变化的“类能量梯度”方向。
真正不同的并不是局部几何本身，而是这种几何如何表示和学习。能量模型参数化标量能量，再通过求导得到向量场；基于得分模型通常直接参数化向量场；扩散模型则更进一步，学习时间依赖的得分场
\[
s_\theta(x,t)
\approx
\nabla_x\log p_t(x),
\]
其中 $p_t(x)$ 是数据经过加噪直到时间 $t$ 后得到的分布。

从能量视角看，每个含噪分布 $p_t$ 都可以与一个隐式能量景观联系起来：
\[
E_t(x)
=
-\log p_t(x)+C_t,
\]
其中 $C_t$ 是任意常数。由于只有梯度起作用，该常数没有影响，于是
\[
\nabla_x\log p_t(x)
=
-\nabla_xE_t(x).
\]
因此，可以把扩散模型解释为学习
\[
s_\theta(x,t)
\approx
-\nabla_xE_t(x),
\]
即一整个由噪声水平索引的能量景观族的负梯度。

这个视角之所以有用，是因为前向加噪过程创造了一套多尺度分布层级。在高噪声水平下，$p_t$ 平滑并接近简单分布，因此对应的能量景观宽缓而规则；在低噪声水平下，$p_t$ 更接近数据分布，因此能量景观包含更精细的结构。扩散训练就是在学习这条序列的局部几何：
\[
\text{粗尺度能量景观}
\quad\longrightarrow\quad
\text{精细的数据级能量景观}.
\]

从这个意义上说，扩散模型可以被看作一种实用且稳定的多尺度能量梯度学习形式。它不是一次性学习一个结构极其尖锐的数据分布能量景观，再直接从中采样；而是把问题分散到许多经过平滑的分布上。高噪声层级提供全局引导，低噪声层级提供细致的局部修正。

这也解释了为什么扩散模型能够绕开传统能量学习中的一些困难。经典 EBM 最大似然需要从当前模型分布采样，以估计负相位；扩散模型则通过已知的前向扰动过程构造含噪训练对，并跨不同噪声层级学习去噪方向。所得方法依然与能量梯度有深刻联系，但训练过程稳定得多。

因此，能量模型为这一部分生成建模提供了更广阔的“景观”视角：它通过标量相容性函数描述分布；而基于得分的扩散，则可以看作现代生成建模中一种尤其有效的方法，用来沿从噪声到数据的多尺度路径，学习并利用相应的能量梯度方向。

\section*{总结与参考文献}
\addcontentsline{toc}{section}{总结与参考文献}

本章超越基于似然的生成建模，介绍了通过比较、几何以及未归一化景观进行学习的模型。GAN 展示了生成器如何通过对抗比较而非似然最大化进行训练。最优判别器把原始 GAN 目标与 Jensen--Shannon 散度联系起来，而实践中通常采用非饱和损失，为生成器提供更强的梯度。

Wasserstein GAN 进一步发展了这一对抗视角，用分布之间的几何距离取代基于重叠的分类。通过对偶形式，可以利用实值的 Lipschitz 评论器估计 $1$-Wasserstein 距离。这个评论器不输出概率；其分数差异反映概率质量如何在数据空间中搬运。

能量模型则提供了更一般的景观视角。EBM 为每个配置赋予标量能量 $E_\theta(x)$，低能量对应合理数据，高能量对应不合理配置。通过 $p_\theta(x)\propto \exp(-E_\theta(x))$，能量定义了未归一化密度。因此，学习与采样围绕能量差异和能量梯度展开，而不是围绕可处理的归一化似然展开。

最后，我们利用恒等式
\[
\nabla_x\log p_\theta(x)
=
-\nabla_xE_\theta(x)
\]
把 EBM 与基于得分的扩散联系起来。这表明能量模型与得分模型共享相同的局部几何：二者都使用把样本引向更高密度、较低能量区域的方向。扩散模型把这一思想扩展到一系列含噪分布，从粗糙的高噪声层级一直学习到精细的数据级结构。

\bigskip

生成对抗网络由 Goodfellow 等人~\cite{goodfellow2014generative} 提出，他们把生成表述为生成器与判别器之间的极小极大博弈，并把最优判别器与 Jensen--Shannon 散度联系起来。Arjovsky 等人~\cite{arjovsky2017wasserstein} 提出了 Wasserstein GAN，以基于最优传输的距离取代基于分类的比较；Gulrajani 等人~\cite{gulrajani2017improved} 随后发展了梯度惩罚形式。

能量建模在机器学习和统计物理中有很长的发展历史。LeCun 等人~\cite{lecun2006tutorial} 从“塑造能量景观”的角度给出了广泛教程。经典例子包括 Hopfield 网络~\cite{hopfield1982neural} 和 Boltzmann 机~\cite{hinton1985learning}。Hyv\"arinen~\cite{hyvarinen2005estimation} 提出得分匹配，用于在不计算配分函数的情况下估计未归一化统计模型；Vincent~\cite{vincent2011connection} 则把得分匹配与去噪联系起来。前面的基于得分扩散章节已经详细讨论了 Song 与 Ermon~\cite{song2019generative} 以及 Song 等人~\cite{song2021scorebased} 的工作；这里再次回顾它们，是因为这些工作把能量梯度思想与多尺度得分学习及扩散模型连接了起来。
